\paragraph{The scalar cost of cancellation between returns.}
Keep the hypotheses and original positive normalizers of
Theorem~\ref{shared-whittaker-scalar-envelope}. The evaluation row of a
seed has norm
\[
 \|\operatorname{ev}_{u,g}\|
 =\mathbf1_{N_fH_{u,f}x_uK}(g_f)\|W_u(g_\infty)\|.
\]
Indeed, in a finite column cell only one dual representation block is
evaluated, and the character and compact transports are unitary.
After an arbitrary actual isometric inclusion
\(J_u:C_u^0\hookrightarrow C_u\), put \(r_u^0=\dim C_u^0\) and,
on \(Y\), define
\[
 w_\delta=\|\operatorname{ev}_{u,\delta y}J_u\|,\quad
 A_u=\sum_\delta w_\delta,\quad B_u=\sum_\delta w_\delta^2,
 \quad m_u=\max_\delta w_\delta,
 \quad H_u=B_u-(2m_u-A_u)_+^2.
\]
All rational cosets \(\delta\in\Delta_u\backslash G(F)\) occur.
Absolute convergence for a finite basis implies \(A_u<\infty\)
almost everywhere; the seed identity gives \(\int_YB_u=c_ur_u^0\).
Take the maximum to be zero when all rows vanish.

\begin{proposition}[The sharp bound from return norms]
\label{shared-seed-return-deficit}
\label{shared-dominant-return-deficit}
With \(K_u=c_u^{-1}J_u^*P_u^*P_uJ_u\), one has
\[
 \operatorname{tr}(I-K_u)_+
 \le h_u:=c_u^{-1}\int_YH_u
 \le\min\left\{r_u^0,
              c_u^{-1}\int_Y(A_u^2-B_u)\right\}.
\]
The new integral is always finite. Replacing the retained row norms by
the full evaluation norms gives a valid, possibly larger bound.
For each finite list of norms, \(H_u\) is the smallest universal
pointwise bound of this kind, when no row directions or phases are known.
\end{proposition}

\begin{proof}
For a nonnegative summable list \(w\), write
\(H(w)=\sum w_j^2-(2\max w_j-\sum w_j)_+^2\).
This function is nondecreasing in each coordinate. If no coordinate
dominates, it equals \(\sum w_j^2\). If one coordinate has value
\(m>s\), where \(s\) is the sum of the others, then
\[
 H(w)=\sum_{j\ne k}w_j^2+2ms-s^2.
\]
The partial derivatives are \(2s\) and
\(2w_j+2(m-s)\), respectively, and the formulas agree at the boundary.
Finite truncation proves the countable case.

For row functionals \(a_j\), let
\(D=\sum a_j^*a_j-(\sum a_j)^*(\sum a_j)\).
For every orthogonal projection \(L\), reverse triangle inequality
and the just-proved monotonicity give
\[
 \operatorname{tr}(LD)
 =\sum_j\|a_jL\|^2-\left\|\sum_j a_jL\right\|^2
 \le H((\|a_jL\|)_j)\le H((\|a_j\|)_j).
\]
Apply this to the evaluation rows and integrate. Both quadratic terms
are integrable by the seed identity and the original \(L^2\) period
hypothesis, so no integrability of the off-diagonal double sum is needed.
Taking the positive spectral projection of \(I-K_u\) proves the claim.
Restriction decreases each row norm, proving the full-column comparison.

Finally \(0\le H\le B\). If \(2m\le A\), then
\(B\le mA\le A^2/2\); otherwise
\((A^2-B)-H=2(s^2-\sum_{j\ne k}w_j^2)\ge0\).
Thus \(H\le A^2-B\). Sharpness holds already for scalar rows:
oppose all other phases to a dominant row; if none dominates, choose
phases closing a polygon. The possible resultant lengths are precisely
\([(2m-A)_+,A]\), by the triangle inequality and induction on the
number of segments. The scalar loss is then exactly \(H\).
\end{proof}

This improves the old pair sum on crowded regions by capping their
contribution at the actual seed mass. For example, one return of value
\(9\) at a point of measure one, and fifty returns of each value
\(1,-1\) at a point of measure \(1/100\), have seed
mass \(82\). The normalized source loss and the new bound are both
\(1/82\); the old bound is \(\min(1,99/82)=1\).

Retain also the actual target projection \(\Pi:H\to H_0\), with
\(n=\dim H_0>0\), \(R_0=\sum_u r_u^0>0\), and diagonal kernels
\(B_0,B_\perp\). Its mass loss is
\[
 L_t=\sum_uc_u^{-1}\|(I-\Pi)P_uJ_u\|_{\mathrm{HS}}^2
       \le\int_YFB_\perp.
\]
For the final maps \(c_u^{-1/2}\Pi P_uJ_u\), at the prescribed
scale \(R_0\), positive-trace subadditivity gives
\[
 g_{R_0}+b_{R_0}\le
 \min\left\{1,\frac{\sum_u h_u+\int_YFB_\perp}{R_0}\right\}
 +\frac1{R_0}\int_Y(F-R_0/n)_+B_0.
\]
In particular, setting
\(\mathcal C_u=c_u^{-1}\int_Y(2m_u-A_u)_+^2\), one has
\(h_u=r_u^0-\mathcal C_u\), and \(\sum_u\mathcal C_u>L_t\)
proves nonvanishing of the final family. These identities use the
retained norms. The arithmetic size of the scalar integrals remains to
be bounded; signed orbital estimates can still be stronger.
